\clearpage
\phantomsection
{\nepalifont
\begin{center}
{\fontsize{16}{18}\selectfont शोधसार}
\end{center}

\vspace{1.5\baselineskip}

साना किसानमा आधारित कृषि प्रणालीमा समयमै बालीको रोग निर्धरण गर्न अझै पनि एक प्रमुख
चुनौती रहेको छ, किनकि विशेषज्ञको सल्लाह भौगोलिक र समयको हिसाबले पहुँचभन्दा बाहिर छ।
यस परियोजनाले वस्तु पहिचान {\normalfont (object detection)}, भिजन-ल्याङ्ग्वेज मोडलिङ
{\normalfont (vision-language modelling)}, ब्याकएन्ड व्यवस्थापन र अफलाइनसमेत चल्ने
मोबाइल प्रविधिलाई नेपालको कृषि सन्दर्भ अनुकूल हुने गरी एउटै बालीको-स्वास्थ्य निर्धरण
प्रणालीमा कसरी एकीकृत गर्न सकिन्छ भन्ने विषयमा अनुसन्धान गर्दछ।

कृषि वैद्यले चारवटा उपप्रणालीहरूलाई {\normalfont image-to-advisory} कार्यप्रवाहमा
जोड्दछ। बोटबिरुवाका ११ विभिन्न भागहरू समेटिएको {\normalfont 30,981} तस्वीर र
{\normalfont 1,47,467} एनोटेसन भएको डेटासेटमा तालिम प्राप्त {\normalfont YOLO}
उपप्रणालीले उत्कृष्ट इपोकमा {\normalfont F1} {\normalfont 0.73} र
{\normalfont mAP@50} {\normalfont 0.76} हासिल गरेको छ। साथै, {\normalfont INT8}
क्वान्टाइजेसनले {\normalfont 1.87} गुणा कम्प्रेसनमा पनि ९८ प्रतिशत पहिचान क्षमता
कायम राख्न सफल भएको छ। {\normalfont Qwen2.5-VL-3B-Instruct} मा {\normalfont LoRA}
एडाप्टर र {\normalfont 4-bit} क्वान्टाइज्ड ट्रेनिङको प्रयोग गरी तयार पारिएको
{\normalfont VLM} उपप्रणालीले {\normalfont 13 GB} {\normalfont GPU} मेमोरी भित्रै
तीन-इपोकको सीमित फाइन-ट्युनिङ पूरा गर्दछ। कम अन्तिम मूल्याङ्कन लस र उच्च टोकन
शुद्धताले यस दिशालाई आशाजनक देखाउँछ। स्वतन्त्र परीक्षण सेटमा देखिएको
{\normalfont 67.7} प्रतिशत पार्स-त्रुटि दरले भने स्ट्रक्चर्ड आउटपुटको विश्वसनीयता
अझै पनि मुख्य चुनौती रहेको पुष्टि गर्दछ। {\normalfont Express/Mongoose} ब्याकएन्डले अपलोड
गरिएका तस्वीरहरूको प्रमाणीकरण गर्दछ, {\normalfont VLM} इन्फरेन्स व्यवस्थापन गर्दछ,
र सुरक्षित {\normalfont API} मार्फत निर्धरण तथा सल्लाह उपलब्ध गराउँछ।
{\normalfont React Native} र {\normalfont Expo} मा निर्मित मोबाइल अनुप्रयोगले
किसानहरूलाई नेपाली र अङ्ग्रेजी दुवै भाषामा सुविधा दिन्छ, जसमा अफलाइन डाटाबेस
{\normalfont (SQLite)}, स्क्यान क्याप्चर, र निर्धरण व्यवस्थित प्रस्तुति समावेश छ।

यसको एकीकृत निर्धरण प्रक्रियाले मोबाइलबाट तस्वीर खिच्नेदेखि ब्याकएन्ड प्रशोधन हुँदै
किसानसम्म नतिजा पुर्याउने कार्यलाई जोड्दछ। समग्रमा यसले {\normalfont YOLO} मोडेल,
{\normalfont VLM} तालिम पाइपलाइन, ब्याकएन्ड र मोबाइल आर्किटेक्चरको सफल कार्यान्वयन
प्रस्तुत गर्दछ। वितरण गरिएको {\normalfont (end-to-end)} एकीकरण र मोबाइल कार्यप्रवाह
प्रत्यक्ष रूपमा परीक्षण गरिएको छ। मुख्य सीमितताहरूको कुरा गर्दा; सीमित तीन-इपोक
{\normalfont VLM} प्रशिक्षणका कारण हालको स्वतन्त्र परीक्षण नतिजा, व्यापक स्वचालित
परीक्षणको आवश्यकता, उत्पादन स्तरको विस्तार, र प्रणालीलाई वास्तविक रूपमा प्रयोगमा
ल्याउनुअघि कृषि विशेषज्ञबाट प्रमाणीकरण अनिवार्य हुनुपर्छ।
}

\vspace{2\baselineskip}
\noindent\textit{Keywords:} Krishi Vaidya, crop disease diagnosis, YOLO, vision-language model, mobile agriculture, backend API
